\documentclass[12pt]{article}

\usepackage[T1]{fontenc}
\usepackage{lmodern}
\usepackage{microtype}
\usepackage[left=1in,right=1in,top=0.88in,bottom=0.88in]{geometry}
\usepackage{setspace}
\usepackage{amsmath,amssymb,mathtools}
\usepackage{mathrsfs}
\usepackage{titlesec}
\usepackage{fancyhdr}
\usepackage[hidelinks]{hyperref}

\setstretch{1.12}
\setlength{\parindent}{0pt}
\setlength{\parskip}{0.72\baselineskip}
\setlength{\abovedisplayskip}{15pt plus 3pt minus 2pt}
\setlength{\belowdisplayskip}{15pt plus 3pt minus 2pt}
\setlength{\abovedisplayshortskip}{14pt plus 2pt minus 2pt}
\setlength{\belowdisplayshortskip}{14pt plus 2pt minus 2pt}
\setlength{\jot}{10pt}
\raggedbottom
\clubpenalty=10000
\widowpenalty=10000
\displaywidowpenalty=10000

\titleformat{\section}
  {\large\bfseries}
  {}{0pt}{}
\titlespacing*{\section}{0pt}{1.6\baselineskip}{0.62\baselineskip}

\pagestyle{fancy}
\fancyhf{}
\renewcommand{\headrulewidth}{0pt}
\fancyfoot[C]{\small\thepage}

\newcommand{\R}{\mathbb R}
\newcommand{\I}{\mathcal I}

\begin{document}

\begin{center}
  {\LARGE\bfseries How the Proof Works}\par
  \vspace{0.4\baselineskip}
  {\normalsize Thomas Birnie}\par
\end{center}

\vspace{0.65\baselineskip}

Fix a smooth, rapidly decreasing, divergence-free initial velocity \(u_0\) and
a viscosity \(\nu>0\). Let \((u,p)\) be the smooth Navier--Stokes solution on
its maximal interval \([0,T_*)\). What matters is what the fluid would have to
do if \(T_*\) were finite.

\section*{The critical height}

Take a vortex as it stretches and thins. Incompressibility forces its core to
narrow as it lengthens. When stretching stays ahead of diffusion, the fluid
turns more rapidly around that shrinking core, so the velocity changes across a
smaller distance. Its gradients can rise even while its kinetic energy remains
finite.

The critical-height problem asks how a Sobolev norm reacts as that concentration
moves to a finer scale. Under the Navier--Stokes rescaling
\[
  u_\lambda(x,t)=\lambda u(\lambda x,\lambda^2t),
  \qquad
  p_\lambda(x,t)=\lambda^2p(\lambda x,\lambda^2t),
\]
nonlinear transport and viscous diffusion both acquire a factor \(\lambda^3\).
Shrinking the scale gives neither one an automatic advantage. The Sobolev norm,
however, changes by
\[
  \|u_\lambda(t)\|_{\dot H^s}
  =
  \lambda^{\,s-1/2}
  \|u(\lambda^2t)\|_{\dot H^s}.
\]
Below \(s=\tfrac12\), concentration makes the norm smaller. Above
\(s=\tfrac12\), it makes the norm larger. At \(s=\tfrac12\), the scale factor
disappears. The norm sees the concentration without weakening or amplifying it
by a power of scale, while nonlinear transport and viscosity remain evenly
matched. This is the critical height. Write
\[
  H(t)
  :=
  \frac12\|\Lambda^{1/2}u(t)\|_2^2,
  \qquad
  \Lambda:=(-\Delta)^{1/2}.
\]

The shrinking core also shortens the available time. Viscous change across a
width \(r\) occurs on the scale
\[
  \tau_\nu(r)=\frac{r^2}{\nu}.
\]
Halving the width quarters that clock. Repeated again and again, this is the
Zeno picture: the fluid would have to pass through smaller widths and steeper
gradients in shorter intervals, all accumulating before a finite time, while
its kinetic energy could remain finite.

In order for a singularity to form, velocity has to blow up. In order for that
to happen, its acceleration would also have to blow up, which means its jerk
has to blow up, and so on, and so on:
\[
  u,
  \qquad
  \partial_tu,
  \qquad
  \partial_t^2u,
  \qquad
  \ldots
\]
As the temporal order rises toward a singularity, we would expect to see less
spatial regularity, not more. Navier--Stokes lets us test that expectation
directly.

\section*{Velocity, acceleration, jerk}

The first derivative is
\[
  \partial_tu
  =
  \nu\Delta u
  -(u\cdot\nabla)u
  -\nabla p.
\]
Differentiate the entire equation once:
\[
  \partial_t^2u
  =
  \nu\Delta\partial_tu
  -(\partial_tu\cdot\nabla)u
  -(u\cdot\nabla)\partial_tu
  -\nabla\partial_tp.
\]
One more derivative gives
\[
\begin{aligned}
  \partial_t^3u
  ={}&
  \nu\Delta\partial_t^2u
  -(\partial_t^2u\cdot\nabla)u
  -2(\partial_tu\cdot\nabla)\partial_tu
  \\
  &-(u\cdot\nabla)\partial_t^2u
  -\nabla\partial_t^2p.
\end{aligned}
\]
By the third line, the pattern is visible. Each time derivative acts on the
entire equation, and the binomial coefficients count every place it can land in
the two velocity factors. The complete pattern is
\[
\begin{aligned}
  \partial_t^{r+1}u
  &+
  \sum_{a=0}^{r}\binom{r}{a}
  \bigl((\partial_t^au)\cdot\nabla\bigr)\partial_t^{r-a}u
  +
  \nabla\partial_t^rp
  \\
  &=
  \nu\Delta\partial_t^ru.
\end{aligned}
\]
Set
\[
  V_r:=\partial_t^ru,
  \qquad
  Q_r:=\partial_t^rp.
\]

\section*{More spatial regularity, not less}

Each new response exposes a spatial Sobolev height. At height \(s\), set
\[
  E_s(t):=\frac12\|\Lambda^su(t)\|_2^2.
\]
Define the nonlinear transfer by
\[
  \mathcal P_s
  :=
  -\operatorname{Re}
  \left\langle
    \Lambda^su,
    \Lambda^s\bigl((u\cdot\nabla)u\bigr)
  \right\rangle_{L^2}.
\]
The viscous term moves up one Sobolev height in this energy:
\[
  \nu\left\langle\Lambda^su,\Lambda^s\Delta u\right\rangle_{L^2}
  =
  -\nu\|\Lambda^{s+1}u\|_2^2
  =
  -2\nu E_{s+1}.
\]
The pressure term drops out of the global pairing because \(\nabla\cdot u=0\).
Pairing the full equation gives
\[
  E_s'=\mathcal P_s-2\nu E_{s+1}.
\]
At critical height, \(H=E_{1/2}\), so
\[
  H'=\mathcal P_{1/2}-2\nu E_{3/2}.
\]
Differentiate this relation and reuse the energy row at each new height. After
\(n\) steps,
\[
  \frac{d^nH}{dt^n}
  =
  (-2\nu)^nE_{n+1/2}
  +
  \sum_{j=0}^{n-1}
  (-2\nu)^j
  \partial_t^{\,n-1-j}\mathcal P_{j+1/2}.
\]
The highest term is the Sobolev energy at height \(n+\tfrac12\). This is the
unexpected reversal: as temporal order rises, Sobolev height rises with it. We
followed velocity, acceleration, jerk, and every later response expecting
spatial regularity to fall away. The equation points upward through the Sobolev
scale.

Apply the calculation to any temporal response. If
\[
  E_{r,s}(t):=\frac12\|\Lambda^sV_r(t)\|_2^2,
\]
and \(\mathcal P_{r,s}\) collects the differentiated transport terms, then
\[
  E_{r,s}'
  =
  \mathcal P_{r,s}
  -2\nu E_{r,s+1}.
\]
The pressure pairing vanishes here as well because \(\nabla\cdot V_r=0\).
The index \(r\) selects the temporal response; \(s\) records its spatial
Sobolev height. Viscosity is the term that links one height to the next in every
row.

That rise matters because, in three dimensions,
\[
  s>k+\frac32
  \quad\Longrightarrow\quad
  H^s(\R^3)\hookrightarrow C^k(\R^3).
\]
For any finite \(k\), choose a height above \(k+\tfrac32\). If all finite
heights reach \(T_*\), then all finite spatial derivatives reach \(T_*\), and
\[
  \bigcap_{m<\infty}H^m(\R^3)
  =
  H^\infty(\R^3)
  \hookrightarrow
  C^\infty(\R^3).
\]
The equations have exposed these heights. The finite rectangles are what keep
each chosen collection controlled through \(T_*\).

\section*{Finite rectangles}

Each pair \((r,m)\) records one concrete statement:
\[
  (r,m)
  \quad\longleftrightarrow\quad
  V_r=\partial_t^ru\in H_x^m.
\]
Temporal order runs in one direction and spatial height in the other. Finite
values \(N\) and \(M\) select a finite block. Each row is carried by an adjacent
trace pair:
\[
  V_n\in L^2(0,T_*;H^{M+1}),
  \qquad
  V_{n+1}=\partial_tV_n\in L^2(0,T_*;H^{M-1})
\]
implies
\[
  V_n\in C([0,T_*];H^M).
\]
One row sits above the desired height, and its time derivative sits below it.
Together they carry the middle trace continuously to the endpoint.

If \(T_*\) were finite, the whole-terminal rectangle theorem would supply these
adjacent pairs for every finite \(N\) and \(M\), with \(0\leq n\leq N\),
across the velocity and pressure rows. That finite collection is
\(\mathcal R_{N,M}[u_0]\).

Two rectangles agree wherever they overlap because both come from the solution
generated by \(u_0\). They assemble into \(\I[u_0]\), whose velocity component
belongs to
\[
  \bigcap_{r,m<\infty}
  H^r(0,T_*;H^m(\R^3)).
\]
Here the temporal responses and spatial Sobolev heights meet. Every finite
choice belongs to one object generated by \(u_0\).

\section*{The smooth endpoint}

Take the zeroth temporal row of the intersection:
\[
  \mathscr S_0:=\pi_0\I[u_0].
\]
This row reaches the endpoint at every finite spatial height:
\[
  u\in C([0,T_*];H^m(\R^3))
  \qquad
  \text{for every finite }m.
\]
The velocity reaches one terminal value
\[
  u_*:=u(T_*)
  \in
  \bigcap_{m<\infty}H^m(\R^3)
  =
  H^\infty(\R^3).
\]
Since \(H^\infty(\R^3)\hookrightarrow C^\infty(\R^3)\), \(u_*\) is smooth.
The singularity picture asked the temporal responses to rise while spatial
regularity disappeared. The completed intersection arrives at \(T_*\) with
every finite spatial height intact.

Local smooth existence from \(u_*\) continues the fluid through
\([T_*,T_*+\delta]\). Uniqueness joins that continuation to the solution on
\([0,T_*]\). The proposed terminal time can be crossed, so a finite \(T_*\)
cannot be maximal:
\[
  T_*=\infty.
\]

\section*{Choose any estimate}

The intersection also contains the familiar finite estimates. Restrict
\(\I[u_0]\) to any finite interval
\([0,T]\), choose a target space \(W^{k,p}\), and choose any finite \(m\) for
which
\[
  H^m(\R^3)\hookrightarrow W^{k,p}(\R^3).
\]
Then
\[
  \I[u_0]
  \longrightarrow
  C_tH_x^m
  \longrightarrow
  L_t^\infty W_x^{k,p}.
\]
Three common choices are
\[
\begin{aligned}
  \I[u_0]
  &\longrightarrow C_tH_x^1
  \longrightarrow L_t^\infty L_x^6,
  \\
  \I[u_0]
  &\longrightarrow C_tH_x^2
  \longrightarrow L_t^\infty L_x^\infty,
  \\
  \I[u_0]
  &\longrightarrow C_tH_x^3
  \longrightarrow L_t^\infty W_x^{1,\infty}.
\end{aligned}
\]

The \(W^{1,\infty}\) choice gives a familiar continuation quantity. Since
\[
  H^3(\R^3)\hookrightarrow W^{1,\infty}(\R^3),
\]
we have
\[
  \|\nabla u(t)\|_{L^\infty}
  \leq
  C\|u(t)\|_{H^3},
\]
and hence
\[
  \int_0^T\|\nabla u(t)\|_{L^\infty}\,dt
  \leq
  CT
  \sup_{0\leq t\leq T}\|u(t)\|_{H^3}
  <\infty.
\]
Any finite classical demand works this way: identify its derivative order and
choose a Sobolev height that embeds into it. The familiar continuation estimates
are read from the intersection afterward: they are derived from \(u_0\) through
the intersection, not the other way around.

The one-line recap is
\[
\begin{aligned}
  u_0
  &\longrightarrow
  \{(V_n,Q_n)\}_{n\geq0}
  \longrightarrow
  \{\mathcal R_{N,M}[u_0]\}_{N,M<\infty}
  \longrightarrow
  \I[u_0]
  \\
  &\longrightarrow
  \mathscr S_0
  \longrightarrow
  u_*\in H^\infty
  \longrightarrow
  T_*=\infty.
\end{aligned}
\]

\end{document}
